% BHU PYQ -- Transcribed & Corrected
% Paper: CHEMJ11 / CHEMN11: Basic Concepts of Chemistry - I (Physical Chemistry)
% Exam: B.Sc. (Hons.) I Semester (NEP-2020) Mid Semester Test, 2024-2025
% Department: Chemistry

\documentclass[11pt,a4paper]{article}
\usepackage[a4paper,top=2.5cm,bottom=2.5cm,left=2.8cm,right=2.8cm,headheight=14pt]{geometry}
\usepackage{amsmath,amssymb}
\usepackage[T1]{fontenc}
\usepackage[utf8]{inputenc}
\usepackage{lmodern,microtype,fancyhdr,enumitem,hyperref,lastpage,parskip}

\hypersetup{
    colorlinks=true,
    urlcolor=black,
    linkcolor=black,
    pdftitle={CHEMN11: Basic Concepts of Chemistry-I (Physical Chemistry) Mid Semester Test -- BHU Sem I 2024-25 (NEP)}
}

\pagestyle{fancy}
\fancyhf{}
\renewcommand{\headrulewidth}{0.4pt}
\renewcommand{\footrulewidth}{0.4pt}
\fancyhead[L]{\small   \textit{Banaras Hindu University}}
\fancyhead[R]{\small   \textit{CHEMN-11: Basic Concepts of Chemistry-I (Mid-Sem 2024-25)}}
\fancyfoot[C]{\small Page  \thepage\ of \pageref{LastPage}}


\newlist{parts}{enumerate}{2}
\setlist[parts,1]{label=(\alph*),leftmargin=*,itemsep=4pt,topsep=2pt}
\setlist[parts,2]{label=(\roman*),leftmargin=*,itemsep=2pt,topsep=2pt}

\newcommand{\pts}[1]{\hfill   \textit{[#1]}}


\begin{document}


\begin{center}
  {\large\bfseries DEPARTMENT OF CHEMISTRY, INSTITUTE OF SCIENCE, BHU}\\[2pt]
  {\normalsize B.Sc. (Hons.) I Semester (NEP-2020) Mid Semester Test: 2024-2025}\\[6pt]
  \rule{0.8\linewidth}{0.5pt}\\[6pt]
  {\large\bfseries Basic Concepts of Chemistry-I (Physical Chemistry)}\\[4pt]
  {\normalsize Date: 20 Dec 2024 \hfill Maximum Marks: 13}
\end{center}

\rule{\linewidth}{0.4pt}

\smallskip



\noindent   \textit{Note: Attempt all the questions. Marks for each question are given in the margin.}

\bigskip




\noindent   \textbf{Question 1.} Complete the following:\pts{3}

\begin{parts}
  \item The expression for the law of corresponding states................
  \item The expression of the Poiseuille equation................
  \item Eotvos's empirical relation between surface tension and temperature................
  \item The expression of collision frequency in terms of pressure....
  \item For the orthorhombic crystal system, axial ratio............ and angles.............
\end{parts}

\medskip\rule{\linewidth}{0.2pt}\medskip




\noindent   \textbf{Question 2.} Derive Graham's law of diffusion using the kinetic gas equation.\pts{1}

\medskip\rule{\linewidth}{0.2pt}\medskip




\noindent   \textbf{Question 3.} What is the excluded volume? Show that the excluded volume is four times the actual volume of a molecule.\pts{2}

\medskip\rule{\linewidth}{0.2pt}\medskip




\noindent   \textbf{Question 4.} How will you determine critical volume experimentally?\pts{2}

\medskip\rule{\linewidth}{0.2pt}\medskip




\noindent   \textbf{Question 5.} Discuss the effect of temperature on viscosity with the help of hole theory.\pts{2}

\medskip\rule{\linewidth}{0.2pt}\medskip




\noindent   \textbf{Question 6.} Explain the steps in converting Weiss indices into Miller indices with a suitable example.\pts{2}

\medskip\rule{\linewidth}{0.2pt}\medskip




\noindent   \textbf{Question 7.} Draw and explain the structure of the diamond.\pts{1}

\vfill
\rule{\linewidth}{0.4pt}

\begin{center}\small
\href{https://bhu-exam.vercel.app/}{Click here to visit our website}\\[4pt]
\href{https://me-aryan.vercel.app/}{Click here to visit our contributor}
\end{center}

\end{document}
